% =============================================================
% §3 问题形式化
% =============================================================
\section{问题形式化}
\label{sec:formulation}

本节将云--边 LLM 推理调度问题形式化。\S\ref{sec:sys-model} 引入部署拓扑与负载的符号；\S\ref{sec:physics-constraints} 将五项 AIGC 物理表达为作用于调度决策的约束与代价函数；\S\ref{sec:sched-problem} 将我们的调度问题表述为多目标优化；\S\ref{sec:mdp} 将其转化为 Markov 决策过程（MDP）以便强化学习求解。表~\ref{tab:notation} 汇总了本节通篇使用的符号。

\begin{table}[t]
  \centering
  \caption{\S\ref{sec:formulation} 所用符号。}
  \label{tab:notation}
  \small
  \setlength{\tabcolsep}{4pt}
  \begin{tabular}{@{}l@{\hspace{8pt}}p{0.60\columnwidth}@{}}
    \toprule
    符号 & 含义 \\
    \midrule
    \multicolumn{2}{@{}l}{\emph{拓扑与服务器}} \\
    $\mathcal{S} = \{s_0, \ldots, s_N\}$
      & 服务器池：云 $s_0$，边缘 $s_1, \ldots, s_N$ \\
    $C_s$, $V_s$, $B_s$
      & $s$ 的算力（TFLOPS）、显存（GB）与带宽（Mbps） \\
    $P^{\text{idle}}_s$, $P^{\max}_s$
      & $s$ 的空闲与峰值功耗（W） \\
    $u_s(t)$, $P_s(t)$
      & $s$ 的利用率与瞬时功率 \\
    $\tau(s, s', x)$
      & 从 $s$ 到 $s'$ 传输 $x$~GB 所需时间 \\
    \addlinespace
    \multicolumn{2}{@{}l}{\emph{模型目录}} \\
    $\mathcal{M} = \{m_1, \ldots, m_K\}$
      & LLM 变体目录 \\
    $W_m$, $C^{\text{cold}}_m$
      & $m$ 的权重体积（GB）与冷加载时间（s） \\
    $\kappa_m$, $\phi_m$, $B^{\max}_m$
      & $m$ 的单 token KV 大小、解码下限、最大 batch \\
    $L_s(t)$
      & 时刻 $t$ 已加载到 $s$ 上的模型集合 \\
    \addlinespace
    \multicolumn{2}{@{}l}{\emph{负载}} \\
    $\mathcal{R}$, $\lambda$
      & 请求流与 Poisson 到达率（req/s） \\
    $r_i = (m_i, L^p_i, L^o_i, a_i)$
      & 目标模型、prompt/输出长度、到达时间 \\
    $t^{\text{pre}}_i$, $t^{\text{dec}}_i$
      & 请求 $r_i$ 的 prefill 与 decode 任务 \\
    $w_t$, $n_t$
      & 任务 $t$ 的工作量（TFLOPS）与 token 数 \\
    \addlinespace
    \multicolumn{2}{@{}l}{\emph{执行与目标}} \\
    $\delta^{\text{cold}}$, $\delta^{\text{kv}}$
      & 冷加载与 KV 迁移延迟 \\
    $b_{s,m,k}(t)$, $\beta_k$
      & 批占用；阶段 $k$ 的开销系数 \\
    $T^{\text{solo}}$, $T^{\text{exec}}_{\text{eff}}$
      & 独占 / 批处理执行时间 \\
    $e^{\text{pre}}_i$; $s^{\text{dec}}_i$, $e^{\text{dec}}_i$
      & prefill 完成时刻；decode 开始与完成时刻 \\
    $\text{TTFT}_i$, $\text{TPOT}_i$
      & 首 token 延迟、每输出 token 时间 \\
    $T^*_{\text{TTFT}}$, $T^*_{\text{TPOT}}$
      & TTFT 与 TPOT 的 \slo{} 阈值 \\
    $\text{SLO}(\pi)$, $E_{\text{tok}}(\pi)$
      & 策略 $\pi$ 下的 \slo{} 达成率与每 token 能耗 \\
    $\pi$; $\mathbf{s}_t$, $a_t$, $r_t$
      & 调度策略；MDP 状态、动作、奖励 \\
    \bottomrule
  \end{tabular}
\end{table}

\subsection{系统模型}
\label{sec:sys-model}

\textbf{部署拓扑。}异构算力池 $\mathcal{S} = \{s_0, s_1, \ldots, s_N\}$，其中 $s_0$ 为云服务器，$s_1, \ldots, s_N$ 为边缘服务器。每台服务器 $s$ 具有固定属性：算力 $C_s$（TFLOPS）、显存容量 $V_s$（GB）、上行带宽 $B_s$（Mbps），以及空闲与峰值功耗 $P^{\text{idle}}_s$、$P^{\max}_s$（W）。网络函数 $\tau(s, s', x)$ 给出在服务器 $s$ 与 $s'$ 之间传输 $x$~GB 数据所需的时间。

\textbf{模型目录。}LLM 变体集合 $\mathcal{M} = \{m_1, \ldots, m_K\}$。每个模型 $m$ 具有权重体积 $W_m$（GB）、冷加载时间 $C^{\text{cold}}_m$（s）、单 token KV cache 大小 $\kappa_m$（MB/token）、单 token 解码下限延迟 $\phi_m$（s/token）以及最大 batch 大小 $B^{\max}_m$。

\textbf{请求负载。}请求流 $\mathcal{R} = \{r_1, r_2, \ldots\}$ 按速率为 $\lambda$ 的 Poisson 过程到达。每个请求 $r_i = (m_i, L^p_i, L^o_i, a_i)$ 包含目标模型 $m_i \in \mathcal{M}$、prompt 长度 $L^p_i$、预期输出长度 $L^o_i$ 与到达时间 $a_i$。目标模型是\emph{请求的一部分}，对应生产 API 中每次调用都指明其模型端点的做法；模型选择发生在上游——调度器只决定每个请求在\emph{哪里}运行，而从不决定由\emph{哪个}模型服务它。

\textbf{任务分解。}每个请求 $r_i$ 分解为两个任务 $(t^{\text{pre}}_i, t^{\text{dec}}_i)$，带有严格依赖边 $t^{\text{pre}}_i \to t^{\text{dec}}_i$；prefill 任务携带 $w^{\text{pre}}_i \propto L^p_i$ 的工作量，decode 任务携带 $w^{\text{dec}}_i \propto L^o_i$ 的工作量，二者引用相同的模型 $m_i$ 与请求 id $i$。整个推理负载因而是若干两两独立的两节点 DAG 构成的森林。

\subsection{AIGC 物理作为约束}
\label{sec:physics-constraints}

第~\ref{sec:intro} 节中的五项 AIGC 物理具体化为以下支配任务在服务器上执行方式的约束与代价项。

\textbf{（P1）模型权重驻留。}服务器 $s$ 在时刻 $t$ 维护一个状态 $L_s(t) \subseteq \mathcal{M}$，表示当前驻留在显存中的模型集合。目标模型为 $m$ 的任务派往 $s$ 时，若 $m \notin L_s(t)$，将支付冷加载惩罚 $C^{\text{cold}}_m$；若加载 $m$ 会超出 $V_s$，仿真器以 LRU 顺序逐出未锁定的模型。形式化的准入延迟为：
\begin{equation}
  \delta^{\text{cold}}_{s,m,t} = C^{\text{cold}}_m \cdot
    \mathbb{1}[m \notin L_s(t)].
\end{equation}

\textbf{（P2）KV cache 本地性。}当 decode 任务 $t^{\text{dec}}_i$ 派往服务器 $s$、而其 prefill $t^{\text{pre}}_i$ 在 $s'$ 上运行时，需支付额外的 KV 迁移开销：
\begin{equation}
  \delta^{\text{kv}}_i = \tau\!\left(s', s,
    L^p_i \cdot \kappa_{m_i} / 1024\right).
\end{equation}
当 $s = s'$ 时，该开销为零（KV cache 已本地）。

\textbf{（P3）连续批处理。}用 $b_{s,m,k}(t)$ 表示在服务器 $s$ 上、对模型 $m$、阶段 $k \in \{\text{pre}, \text{dec}\}$ 当前运行中的同模型同类任务数。时刻 $t$ 准入到 $(s, m, k)$ 的新任务经历的有效执行时间为
\begin{equation}
  T^{\text{exec}}_{\text{eff}}
    = T^{\text{solo}} \cdot
      \bigl(1 + b_{s,m,k}(t) \cdot \beta_k\bigr),
\end{equation}
其中 $\beta_{\text{pre}} = 0.30$、$\beta_{\text{dec}} = 0.05$ 为阶段特异的开销系数，按 vLLM 测量校准（\S\ref{sec:simulator}）。准入要求 $b_{s,m,k}(t) < B^{\max}_m$（批槽未满）。

\textbf{（P4）显存带宽下限。}单阶段任务的"独占"执行时间取算力受限时间与带宽下限两者之最大：
\begin{equation}
  T^{\text{solo}}
    = \max\!\left(\frac{w_t}{C_s},\;
                  \phi_{m_t} \cdot n_t\right),
\end{equation}
其中 $w_t$ 为任务工作量（TFLOPS），$n_t$ 为 token 数（prefill 为输入、decode 为输出）。

\textbf{（P5）两阶段依赖。}decode 任务 $t^{\text{dec}}_i$ 只有在 $t^{\text{pre}}_i$ 完全完成后方可 READY。结合（P1）与（P2），派往 $s \neq s'$ 的 decode 无法享受 prefill 的 KV cache——要么支付迁移开销，要么重做 prefill。

\subsection{调度问题}
\label{sec:sched-problem}

调度策略 $\pi$ 为每个 READY 任务在每个时间步给出分派 $\pi: t \mapsto s \in \mathcal{S}$。给定请求流 $\mathcal{R}$，$\pi$ 诱导出：

\begin{itemize}
\item 每请求 \textbf{TTFT}（首 token 延迟）：
  $\text{TTFT}_i = e^{\text{pre}}_i - a_i$，其中
  $e^{\text{pre}}_i$ 为 prefill 任务的实际完成时刻。

\item 每请求 \textbf{TPOT}（每输出 token 时间）：
  $\text{TPOT}_i = (e^{\text{dec}}_i - s^{\text{dec}}_i) / L^o_i$。

\item \textbf{\slo{} 达成率}：
  \begin{align}
    \text{SLO}(\pi) = \frac{1}{|\mathcal{R}|} \sum_i
      \mathbb{1}\!\bigl[&\text{TTFT}_i \le T^*_{\text{TTFT}}
      \notag\\
      &\land \text{TPOT}_i \le T^*_{\text{TPOT}}\bigr].
  \end{align}

\item \textbf{每 token 能耗}：
  \begin{equation}
    E_{\text{tok}}(\pi) =
      \frac{\sum_s \int_0^T P_s(t)\,dt}{\sum_i L^o_i},
  \end{equation}
  其中瞬时功率
  $P_s(t) = P^{\text{idle}}_s + (P^{\max}_s - P^{\text{idle}}_s)
  \cdot u_s(t)$ 依赖于各服务器算力利用率 $u_s(t)$。
\end{itemize}

我们的调度问题是\textbf{双目标规划}：
\begin{equation}
  \max_\pi\; \text{SLO}(\pi)
  \qquad
  \min_\pi\; E_{\text{tok}}(\pi),
\end{equation}
受每步准入可行性约束（算力、显存、批槽）。由于两个目标可能冲突（如"全派往云端"最大化单请求算力但能耗超支），我们在 $(\text{SLO}, E_{\text{tok}})$ 平面上评估策略，并追求\textbf{Pareto 占优}：$\pi$ 占优 $\pi'$ 当且仅当 $\pi$ 在两个轴上均不劣、至少在一个轴上严格优。这两个目标将 \S\ref{sec:intro} 中的延迟与成本动机操作化；隐私需求（若存在）则作为可用服务器集合上的可行性约束进入问题，而非作为一个目标。

\subsection{用于 RL 求解的 MDP 形式化}
\label{sec:mdp}

我们将单任务分派视作序贯决策过程。在每个 READY 任务到达时刻，agent 观测状态 $\mathbf{s}_t$——概括当前任务属性以及所有服务器的利用率、已加载模型集合与批占用（\S\ref{sec:rl}）——输出动作 $a_t \in \mathcal{S}$ 选定一台服务器，并获得标量奖励 $r_t$：7 项物理与 AIGC 感知信号的加权和（\S\ref{sec:rl}），提供\textbf{每决策的稠密}反馈，而非 episode 末的 \slo{}/能耗结果——考虑到 episode 视界很长（每次运行约 100 次分派），这一点至关重要。

该 MDP 在原则上\textbf{部分可观}，但可观特征已涵盖仿真器动力学下全部决策相关信息；它在 episode 内非平稳（利用率随时间演化）、在 episode 间平稳（负载分布固定）。我们用带动作掩码、预训练与广义优势估计（GAE）的 PPO 求解，详见 \S\ref{sec:rl}。

\textbf{为什么选 RL。}RL 的一个自然替代是直接求解 \S\ref{sec:sched-problem} 隐含的整数线性规划（ILP）。我们放弃该方案的原因：（i）AIGC 物理（P1）（P3）引入\emph{状态相关}代价（冷加载与批处理开销依赖服务器当前运行集合，且在 episode 内不断变化），破坏了静态代价 ILP 的假设；（ii）每步的组合动作空间（$|\mathcal{S}|^{|\mathcal{R}|}$）在真实负载规模下对分支定界而言过大。
